\documentclass[a4paper,fleqn]{cas-dc}
\usepackage[authoryear,longnamesfirst]{natbib}
\begin{document}
\let\WriteBookmarks\relax
\shorttitle{Glacier retreat in a fictional valley}
\shortauthors{K. Fujimori et al.}
\title [mode = title]{Two decades of glacier retreat and proglacial lake growth in a fictional mountain valley}
\author[1]{Kazuha Fujimori}[orcid=0000-0000-0000-0001]
\cormark[1]
\ead{fujimori@example.org}
\author[2]{Anders Nyg{\aa}rd}
\author[1]{Mei Takeyama}
\affiliation[1]{organization={Graduate School of Earth Science, Kitamachi University},
            city={Kitamachi},
            postcode={100-0000},
            country={Japan}}
\affiliation[2]{organization={Institute of Fictional Glaciology, University of Nowhere},
            city={Nowhere},
            postcode={00000},
            country={Freedonia}}
\cortext[1]{Corresponding author}
\begin{abstract}
Mountain glaciers are retreating worldwide, and the lakes that form in front of them pose a growing flood hazard.
We combine two decades of fictional satellite imagery with field surveys to quantify the retreat of five glaciers in the imaginary Kitamachi Valley.
Between 2004 and 2024 the glaciers lost 31\% of their area, while the largest proglacial lake tripled in volume.
Retreat accelerated after 2015 and was strongly correlated with summer air temperature at a nearby station.
A simple hydraulic model indicates that a moraine breach at the largest lake could produce a peak discharge ten times higher than the annual flood.
We recommend continuous monitoring of the lake level and of the moraine dam.
\end{abstract}
\begin{keywords}
glacier retreat \sep proglacial lakes \sep remote sensing \sep natural hazards
\end{keywords}
\maketitle
\section{Introduction}
Glaciers are sensitive indicators of climate change, and their retreat affects water resources and natural hazards in mountain regions \citep{fujimori2021,nygaard2022}.
As glaciers shrink, meltwater often accumulates behind unstable moraines, forming lakes that may drain suddenly.
Such glacial lake outburst floods have caused severe damage in several mountain ranges \citep{takeyama2023}.

Here we document the evolution of glaciers and lakes in a fictional valley over two decades.
We use optical imagery to map glacier outlines and lake areas, bathymetric surveys to estimate lake volumes, and a hydraulic model to assess the flood hazard.

\section{Study area}
The Kitamachi Valley lies at an elevation of 3000--5000~m.
\bibliographystyle{cas-model2-names}
\bibliography{refs}
\end{document}
